% ============================================================
%  METHODS  (revised)
%  Drop-in replacement for the methods section.
%  Requires: amsmath, booktabs, tikz, tikzlibraries {arrows.meta, positioning, shapes.geometric}
% ============================================================

\section{Methods}

This paper aims to explore how heterogeneous information sources and dynamic trust
relationships influence collective decision-making in disaster response.
To this end, an agent-based model is built that combines opinion dynamics frameworks
with a reinforcement-learning decision layer that captures the opportunities and
challenges of human-AI interaction in crises.
The model simulates decentralised disaster relief efforts where human agents,
characterised by distinct behavioural types, interact with human and AI information
sources to allocate emergency resources across a spatial environment.

The agent-based modelling (ABM) approach allows us to analyse the trade-offs between
AI adoption, resulting situational awareness, filter-bubble formation, and relief
performance.
By identifying the \emph{Goldilocks zone} where AI alignment transitions from
amplifying to mitigating filter bubbles, we contribute insights to ongoing discussions
on trustworthy AI and objective trade-offs in morally charged crisis-management
scenarios.
The model expands established opinion dynamics frameworks \citep{geschke2019triple}
to incorporate the challenges of human-AI interaction in crisis contexts.

% ============================================================
\section{Modelling Framework}

The simulation is situated on a $30\times30$ grid where disaster severity varies from
level~0 (no damage) to level~5 (critical need).
Human agents of two behavioural types---exploitative and exploratory---navigate this
environment alongside AI agents that provide variable-quality information.
The model proceeds in discrete time steps where agents sequentially sense their
immediate environment, seek information from social or AI sources, and allocate
relief tokens to perceived high-need locations.

The choice of Q-learning for source selection and precision-weighted Bayesian updating
for belief revision represents a novel integration of reinforcement learning and
probabilistic reasoning in opinion dynamics modelling.
While traditional opinion dynamics models such as \citet{deffuant2000mixing} and
\citet{hegselmann2002opinion} typically assume simple averaging or binary bounded
confidence, our approach captures the strategic dimension of information seeking and
the uncertainty inherent in emergency decision-making.
This methodology bridges classical models in social physics with contemporary
frameworks for modelling human-AI interaction.

% ============================================================
\subsection{Agent Architecture}

Human agents maintain cognitive representations comprising individual belief maps
about disaster severity across all grid locations and trust profiles for each
potential information source.
Each agent operates with a Q-learning table encoding the expected value of
consulting different source modes (AI, social contact, self-action), while trust
values reflect accumulated assessments of source accuracy and reliability.

Although both exploitative and exploratory agents share the same direct-sensing
radius (own cell only), they differ in \emph{interest-point targeting}: exploitative
agents focus their information-seeking on the \emph{believed epicentre}---the
cell they currently judge most critical---while exploratory agents direct queries
towards the \emph{highest-uncertainty area}, the cell whose belief has the greatest
unresolved variance.
This distinction, rather than a difference in raw sensing range, produces the
divergence in information strategies: exploitative agents confirm and refine a focal
assessment, while exploratory agents actively reduce epistemic uncertainty.

The two types also differ in their acceptance parameters (Table~\ref{tab:S2}).
Exploitative agents have a narrow acceptance window ($D=2.0$) and steep
sensitivity ($\delta=3.5$), making them resistant to information that departs
substantially from their prior.
Exploratory agents have a wide, gradual acceptance curve ($D=4.0$, $\delta=1.2$)
that readily incorporates novel reports.
These archetypes map onto documented individual differences in information
processing under uncertainty \citep{HertwigEngel2016,March1991}.

AI agents complement the system by sensing 15\% of the grid per tick and responding
to human queries with reports whose truth content is controlled by the alignment
parameter $\alpha$ (Section~\ref{sec:alignment}).

% ============================================================
\subsection{Information Diffusion Mechanisms}
\label{sec:alignment}

The alignment parameter $\alpha\in[0,1]$ governs the truthfulness of AI reports.
A queried AI agent responds with
%
\begin{equation}
  r_{\mathrm{AI}} = (1-\alpha)\,t + \alpha\,b,
  \label{eq:alignment}
\end{equation}
%
where $t$ is the cell's true severity and $b$ is the querying agent's current
belief.
At $\alpha=0$ the AI reports ground truth; at $\alpha=1$ it perfectly confirms the
agent's prior; intermediate values interpolate linearly.
This formulation isolates the confirmation effect from other AI failure modes
(latency, coverage, bias), enabling a clean sweep across the truth--confirmation
spectrum.

\begin{figure}[h]
\centering
\begin{tikzpicture}[
    node distance=1.1cm and 2cm,
    every node/.style={font=\sffamily\small, align=center},
    process/.style={rectangle, draw, fill=blue!10, text width=4.2cm,
                    minimum height=1cm, rounded corners},
    decision/.style={diamond, draw, fill=orange!10, text width=2.2cm,
                     inner sep=0pt, aspect=2},
    arrow/.style={-{Stealth[scale=1.2]}, thick}
]

\node (start)  [circle, draw, fill=gray!20]                             {Start};
\node (step1)  [process, below=of start]
    {\textbf{1. Sense Environment}\\Radius: 0 (own cell only)\\Noise: 8\%};
\node (step2)  [decision, below=of step1, yshift=-0.4cm]
    {\textbf{2. Select Source}\\ $\varepsilon=0.30$};
\node (step3)  [process, below=of step2, yshift=-0.4cm]
    {\textbf{3. Query Source}\\(Self\,/\,Friend\,/\,AI)\\AI: Eq.\,\eqref{eq:alignment}};
\node (step4)  [process, below=of step3]
    {\textbf{4. Update Beliefs}\\Bayesian, precision-weighted\\Eq.\,\eqref{eq:acceptance}--\eqref{eq:bayesian}};
\node (step5)  [process, below=of step4]
    {\textbf{5. Send Relief}\\Top 5 cells ($L\!\geq\!3$)\\Queue reward};
\node (step6)  [process, below=of step5]
    {\textbf{6. Process Rewards}\\15--25 tick delay\\Update $Q$ \& Trust};
\node (step7)  [process, below=of step6]
    {\textbf{7. Apply Decay}\\Trust: type-specific rate\\$(0.015$ / $0.030)$};

\draw [arrow] (start) -- (step1);
\draw [arrow] (step1) -- (step2);
\draw [arrow] (step2) -- (step3);
\draw [arrow] (step3) -- (step4);
\draw [arrow] (step4) -- (step5);
\draw [arrow] (step5) -- (step6);
\draw [arrow] (step6) -- (step7);
\draw [arrow] (step7.west) -- ++(-1.6,0) |- (step1.west)
    node[pos=0.25, left, rotate=90] {Next tick};
\end{tikzpicture}
\caption{Per-tick decision cycle for human agents.
         Source selection uses an $\varepsilon$-greedy Q-learning strategy.
         Belief updating follows the precision-weighted Bayesian mechanism
         (Equations~\ref{eq:acceptance}--\ref{eq:bayesian}).
         Reward processing is delayed 15--25 ticks to model logistical
         pipelines; trust updates use agent-type-specific learning rates.}
\label{fig:cycle}
\end{figure}

When agents receive a report $r$ from any source, they accept it with probability
%
\begin{equation}
  P(\text{accept}) = \frac{D_{\mathrm{eff}}^{\delta}}
                          {|r-b|^{\delta}+D_{\mathrm{eff}}^{\delta}},
  \label{eq:acceptance}
\end{equation}
%
where $D_{\mathrm{eff}} = D\cdot(1+0.5\,T_{\mathrm{source}})$ for social contacts
(trust-widened) and $D_{\mathrm{eff}}=D$ for the AI.
If the report is accepted, the posterior is precision-weighted:
%
\begin{equation}
  b' = \frac{\pi_{\mathrm{prior}}\,b + \pi_{\mathrm{source}}\,r}
            {\pi_{\mathrm{prior}} + \pi_{\mathrm{source}}},
  \label{eq:bayesian}
\end{equation}
%
with $\pi\propto c/(1-c)$ (confidence odds), scaled by a type-specific factor
(1.5 for exploitative, 0.8 for exploratory) that calibrates resistance to belief
revision.
This mechanism generalises the Deffuant bounded-confidence model
\citep{deffuant2000mixing} by replacing binary thresholding with a smooth
probability curve and by coupling the effective threshold to dynamic source trust,
producing richer dynamics while remaining analytically interpretable
\citep{lorenz2017continuous}.

When selecting information sources, agents employ an $\varepsilon$-greedy strategy
($\varepsilon=0.30$) that balances exploration of unfamiliar sources with
exploitation of trusted relationships.
The social network is constructed with type-homogeneous community structure:
exploitative agents are wired predominantly within their own cluster, exploratory
agents within theirs, with no enforced cross-type edges.
This follows empirical evidence of homophily in disaster communication networks
\citep{Bruns2012,Sutton2008}.

% ============================================================
\subsection{Trust Evolution}

Trust dynamics emerge from two parallel processes: direct observation of
information accuracy and reinforcement learning based on relief outcomes.
The trust update mechanism combines immediate feedback from ground-truth comparisons
with delayed rewards from the effectiveness of relief allocation.
This dual-track approach captures both the cognitive updating of source reliability
assessments and the pragmatic evaluation of information utility for decision support.

The model uses asymmetric trust learning rates across agent types:
exploitative agents update trust slowly ($\mathrm{lr}=0.015$), consistent with
their confirmation-seeking profile; exploratory agents update more rapidly
($\mathrm{lr}=0.030$), reflecting their higher sensitivity to new evidence.
This behavioural heterogeneity produces emergent trust differentiation that
significantly influences collective performance under varying alignment conditions.

% ============================================================
\subsection{Echo Chamber Quantification}

The Social Echo Chamber Index (SECI) measures within-community belief homogeneity
relative to population-wide variance, following the approach of
\citet{cinelli2021echo}.
Because the social network is type-homogeneous, SECI is computed separately for
each connected component (community) and then averaged by agent type.
For a community with internal belief variance $V_c$ and global population variance
$V_g$:
%
\begin{equation}
  \mathrm{SECI} =
  \begin{cases}
    \displaystyle\max\!\left(-1,\;\frac{V_c - V_g}{V_g}\right)
      & V_c < V_g \quad\text{(echo chamber)}, \\[8pt]
    \displaystyle\min\!\left(+1,\;\frac{V_c - V_g}{V_{\max} - V_g}\right)
      & V_c \geq V_g \quad\text{(anti-echo)},
  \end{cases}
  \label{eq:seci}
\end{equation}
%
where $V_{\max}=5.0$ is the maximum possible variance on the $[0,5]$ severity
scale.
$\mathrm{SECI}<0$ indicates community beliefs are more homogeneous than the global
pool (echo chamber); $\mathrm{SECI}=0$ is the null; $\mathrm{SECI}>0$ indicates
the community is more internally diverse (anti-echo).
The piecewise normalisation keeps the index symmetric on $[-1,+1]$ without
assuming that the maximum possible variance equals the global variance.

Alternative formulations based on opinion fragmentation \citep{Bail2018} or
information entropy \citep{Sasahara2021} were considered; the variance-based
approach was preferred because it has a natural zero, is directly comparable across
agent types, and maps onto the precision terms in the Bayesian update
(Equation~\ref{eq:bayesian}).

The AI Echo Chamber Index (AECI) measures the fraction of information updates
that are actually absorbed from AI sources within each five-tick window:
%
\begin{equation}
  \mathrm{AECI} = \frac{\Delta q_{\mathrm{AI}}}
                       {\Delta q_{\mathrm{AI}} + \Delta q_{\mathrm{human}}},
  \label{eq:aeci}
\end{equation}
%
where $\Delta q$ counts D$\delta$-accepted updates during the window.
Unlike a simple query-count ratio, this formulation distinguishes
\emph{reliance}---AI information that passes the acceptance threshold and
revises beliefs---from mere querying, which may be rejected without effect.
High AECI signals over-dependence on a potentially confirming source, regardless
of how often human contacts are also consulted.

% ============================================================
\subsection{Relief Allocation and Performance Evaluation}

Agents allocate relief tokens to the top five cells they judge to have severity
$L\geq3$.
Exploitative agents score cells by confidence-weighted severity; exploratory agents
add an uncertainty-reduction term that rewards directing resources to cells with
high epistemic ambiguity.
Token placement effectiveness is evaluated through delayed reward signals
(15--25 ticks) that compare allocated resources against actual need distributions
at the time of delivery.

Performance metrics encompass belief accuracy (mean absolute error against ground
truth), targeting precision (fraction of tokens directed to cells with actual
severity $\geq3$), unmet critical needs (cells with severity $\geq4$ receiving zero
tokens per tick), and spatial coverage deficit (per-cell shortfall averaged near
and far from the disaster epicentre).
Together these measures capture both efficiency and equity dimensions of disaster
response, crucial for evaluating the downstream consequences of information flow.

% ============================================================
\subsection{Experimental Framework}

The study implements two core experiments and a set of sensitivity sweeps.

\paragraph{Goldilocks alignment sweep.}
The primary experiment sweeps $\alpha$ across eleven levels ($0.0,0.1,\ldots,1.0$)
with $N=20$ independent replications per level.
Replications differ in epicentre location, agent initialisation, and stochastic
shock sequences.
Optimal alignment is identified by two composite scores, both constructed over the
steady-state window (last 75 of 200 ticks) using range normalisation across the
sweep so each component contributes on a $[0,1]$ scale:
%
\begin{align}
  \text{total\_bubble}(\alpha)
    &= |\mathrm{SECI}|_{\mathrm{norm}}(\alpha)
       + |\mathrm{AECI}|_{\mathrm{norm}}(\alpha), \label{eq:bubble}\\
  \text{total\_score}(\alpha)
    &= |\mathrm{SECI}|_{\mathrm{norm}}(\alpha)
       + |\mathrm{AECI}|_{\mathrm{norm}}(\alpha)
       + \mathrm{MAE}_{\mathrm{norm}}(\alpha). \label{eq:score}
\end{align}
%
The Goldilocks optimum $\alpha^*$ minimises total\_bubble; $\alpha^*(+\mathrm{MAE})$
minimises total\_score.
Reporting both criteria makes the trade-off between echo-chamber suppression and
belief accuracy explicit.

\paragraph{Cognitive gap sweep.}
A secondary experiment varies the cognitive gap scalar $g\in\{0.0,0.5,1.0,1.5\}$,
which scales the divergence between exploitative and exploratory acceptance
parameters from a shared midpoint while preserving their ordinal relationship
($D_{\mathrm{exploit}}<D_{\mathrm{explor}}$,
$\delta_{\mathrm{exploit}}>\delta_{\mathrm{explor}}$).
At $g=0$ both types are cognitively identical, providing a null condition for
heterogeneity effects; $g=1$ recovers the baseline calibration
(Table~\ref{tab:S4}).
This sweep tests whether $\alpha^*$ is robust to within-population cognitive
diversity.

\paragraph{Factor sensitivity sweeps.}
One-at-a-time sweeps over agent-type composition, disaster dynamics, and rumour
probability (detailed in Table~\ref{tab:S8}) test the generalisability of the
Goldilocks finding across environmental and population variation.

Statistical analysis employs non-parametric methods to handle non-normal
performance distributions across the 20 replications, while time-series analysis
tracks the evolution of trust and bubble metrics throughout each crisis scenario.
Spatial and network periphery analyses decompose outcomes by distance to the
disaster epicentre (nearest vs.\ furthest quartile of cells) and by social-network
degree (Q1 vs.\ Q4 agents) to examine whether alignment benefits are equitably
distributed.

% ============================================================
\subsection{Technical Implementation}

The model is implemented in Python~3.11 using the Mesa~3.x framework for agent
scheduling, NetworkX for community detection and social-network analysis, and
NumPy / Matplotlib for computation and visualisation.
All simulation runs are seeded via \texttt{numpy.random.seed(run\_index)},
ensuring bit-for-bit reproducibility.
Results are serialised to JSON after each $(\alpha, \text{run})$ cell; the
plotting pipeline reads these files without re-running the simulation, so all
figures can be regenerated without recomputation.
Full parameter tables and software dependencies are provided in the supplementary
materials (Tables~S1--S2, S9).
